\documentclass{ximera}

\newcommand{\inbetween}{\fontsize{7pt}{9pt}\selectfont}


\ifdefined\HCode
  \newcommand{\alttext}[1]{\HCode{<span class="sr-only">#1</span>}}
\else
  \newcommand{\alttext}[1]{} % Fallback for PDF view
\fi

% MATH -----------------------------------------------------------
\newcommand{\nats}{\mathbb N}
\newcommand{\ints}{\mathbb Z}
\newcommand{\rats}{\mathbb Q}
\newcommand{\reals}{\mathbb R}
\newcommand{\complex}{\mathbb C}
\newcommand{\powerset}{\mathscr P}

%------------------------------------------------------------



\newcommand{\Laplace}[1]{\ensuremath{\mathcal{L}{\left[#1\right]}}}
\newcommand{\InvLap}[1]{\ensuremath{\mathcal{L}^{-1}{\left[#1\right]}}}


\pgfplotsset{compat=1.18}

\begin{document}
\begin{abstract}
 \end{abstract}

    \title{Lab 7 - Convolution}
    
    \maketitle


\begin{enumerate}
    
%\item Determine the displacement above equilibrium $y(t)$ (at $t\geq 0$ seconds) of a $0.1$ kg mass in an undamped spring-mass system with a natural frequency of $\displaystyle \frac{5}{\pi}$ hz, if the mass begins at rest at equilibrium and is forced (in Newtons) by $\displaystyle F(t)=\left\{\begin{array}{lr}
%       t & \mbox{ if }0\leq t<1\\
%       \\
%       1 & \mbox{ if }1\leq t<2\\
%       \\
%       2-\dfrac{t}{2} &  \mbox{ if }2\leq t<4
%       \\
%       0  & \mbox{ if }t\geq 4\\
%     \end{array}\right.$  


%\begin{figure}[h]
%\alttext{The graph of a piecewise continuous function.}
% \centering
% \begin{tikzpicture}[scale=0.75]
%\begin{axis}[axis x line=middle, axis y line=
%middle, xtick={-1,1,2,3,4,5,6}, ytick={-1,...,2,3},
%xmin=-2, xmax=7, ymin=-0.5, ymax=1.5]

%\addplot[blue,thick,domain=0:1] {x};
%\addplot[blue,thick,domain=1:2] {1};
%\addplot[blue,smooth,thick,domain=2:4] {2-0.5*x};
%\addplot[blue,smooth,thick,domain=4:7] {0};
%\end{axis}
%\end{tikzpicture}
%\caption{Graph of a piecewise continuous function}
%\end{figure}


%\vfill

%\item Solve the IVP $y^{\,\prime\prime}+2y^{\,\prime}+y=e^{-t}+\delta((t-1)$.

%\vfill

%\newpage

\item Let $f(t),g(t)$ be functions.  The \textbf{convolution} of $f$ and $g$ is the function $$\displaystyle (f\ast g)(t)=\int_{0}^{t}f(x)g(t-x)\,\mathrm{d}x.$$

\begin{enumerate}

\item  Compute the convolution of $f(t)=t^2$ and $g(t)=e^{-t}$ directly from the definition above.  Put the answer in the box provided.

\hfill{}\fbox{ \begin{minipage}{4in}\vspace{0.3in} \hfill\vspace{0.3in} \end{minipage} }

\vfill

\item Show that the convolution of $f(t)$ and $g(t)$ does not depend on order, that is, $f\ast g=g\ast f$.  

Hint:  use the $u$-substitution $u=t-x$.

\vfill


\end{enumerate}


\newpage


Recall: the \textbf{Laplace transform} of a function $f(t)$ is $\displaystyle \mathcal{L}[f]=\int_{0}^{\infty}f(t)e^{-st}\,\mathrm{d}t$.

\item  Let $f(t)$ and $g(t)$ be continuous functions for $t\geq 0$.  Fill-in the missing details below, after each of the ellipses, which will demonstrate that\\ $\mathcal{L}[f\ast g]=\mathcal{L}[f]\mathcal{L}[g]$:


First note that $\displaystyle \mathcal{L}[f]\mathcal{L}[g]=\int_0^{\infty}f(t)e^{-st}\,\mathrm{d}t\int_0^{\infty}g(t)e^{-st}\,\mathrm{d}t$.  By replacing all instances of $``$t" with $``$x" in the first integral and $``$y" in the second integral, we get that


\vfill

$\displaystyle \mathcal{L}[f]\mathcal{L}[g]=$ ... \hfill.

\vfill

Bringing the second integral into the integrand of the first integral, 

\vfill

$\displaystyle \mathcal{L}[f]\mathcal{L}[g]=$ ...\hfill.  

\vfill


Then by bringing the integrand of the first integral into the second integral, we get the iterated double integral 

\vfill

$\displaystyle \mathcal{L}[f]\mathcal{L}[g]=\int_0^{\infty}\int_0^{\infty}$ ...\hfill $\mathrm{d}y\,\mathrm{d}x$.  

\vfill


Using the substitution $t=x+y$ in the $``$inner" integral, ($x$ can be regarded as a constant).  


\vfill

Then $\mathrm{d}t=\mathrm{d}y$, and $t=$ ... \hfill when $y=0$, and $t\rightarrow ...$\hfill when $y\rightarrow\infty$, 

\vfill

so we get $\displaystyle \mathcal{L}[f]\mathcal{L}[g]=\int_0^{\infty}\int_x^{\infty}$ ...\hfill $\mathrm{d}t\,\mathrm{d}x$.  

\vfill


The region of the $xt$-plane for this improper iterated double integral is ... \\(\emph{Shade the region!})

\begin{figure}[h]
\alttext{The identity function through the first quadrant of a plane.}
 \centering
 \begin{tikzpicture}[scale=0.55]
\begin{axis}[
axis x line=middle,
axis y line=middle,
ylabel=$x$,
xlabel=$t$,
xtick={0},
ytick={0},
xmax=10,
ymax=10,
xmin=-1,
ymin=-1
]
% Plots
\addplot[blue,ultra thick,domain=0:9] {x};
\end{axis}
\end{tikzpicture}
\caption{A line in a a plane}
\end{figure}

\newpage

Let's reverse the order of the integration, from $\mathrm{d}t\,\mathrm{d}x$ to $\mathrm{d}x\,\mathrm{d}t$.  The region shaded can be described as 

\vfill


$\{(t,x)|0\leq t\leq \infty,$ ...\hfill$\leq x\leq$ ...\hfill $\}$ so we get that


\vfill

 $\displaystyle \mathcal{L}[f]\mathcal{L}[g]=$ ...  \hfill 
 
\vfill

 
 and the inner integral is not improper!  Pulling $e^{-st}$ out of the inner integral and using $\displaystyle (f\ast g)(t)=\int_{0}^{t}f(x)g(t-x)\,\mathrm{d}x$ we get 
 
\vfill

 $\displaystyle \mathcal{L}[f]\mathcal{L}[g]=$ ...  \hfill  $=\mathcal{L}[...$\hspace{3cm} $]$.


\vfill


\item Now compute the convolution of $f(t)=t^2$ and $g(t)=e^{-t}$ as follows:

\begin{itemize}
\item Determine the product $\mathcal{L}(f)\mathcal{L}(g)$.

\item Perform a partial fraction decomposition on the above function.

\item Take the inverse Laplace transform of the result.

\end{itemize}

\vfill
\vfill

\item Solve the ``integral equation" $y(t)=2+\displaystyle \int_0^t (t-x)y(x)\,\mathrm{d}x$.  Put the answer in the box provided.

\hfill{}\fbox{ \begin{minipage}{4in}\vspace{0.3in} \hfill\vspace{0.3in} \end{minipage} }


Hint: the integral is the convolution of two functions.

\vfill
\vfill


\end{enumerate}




\end{document}